

\Bsec{Fourier-Transformation}{
		
		\begin{itemize}
			\item Wdh. Komplexe Zahlen
			\item DFT vs. IDFT
			\item TBD: Schnelle inverse Fourier Transformation
			\item Verständnis/Trick Fragen
		\end{itemize}
		}{}
		
%----------------------------------------------------------------------------------
\bsubsec{Wdh. Komplexe Zahlen}{
		
		$$z = x + iy$$
		$$x = Re(z); y = Im(z)$$
		$$i = \sqrt{-1}$$
		$$\compl{z} = z - iy$$
		Euler Formel: $$e^{i \cdot m} = cos(m) + i \cdot sin(m)$$
		
	}
	\btask{ 
		
		\begin{enumerate}
			\item Berechne $(2 + i3) + (1 - i2)$. Wie groß ist der Imaginärteil?
			\item Was ist $e^{i \cdot 2\pi} * 2$
			\item Berechne den Betrag von $4 + i3$!
		\end{enumerate}
		
	}
	\bsol{
		
		\begin{enumerate}
			\item $= 3 + i$. $Im(3 + i) = 1$ (und nicht $1i$!!)
			\item $= 2$ (Imaginärteil wird zu 0)
			\item $\abs{4 + i3} = 5$
		\end{enumerate}
	}
%----------------------------------------------------------------------------------
\bsubsec{DFT vs. IDFT}{
		
		Diskrete Fourier Transformation:
		$$\begin{bmatrix}
		c_0 \\ c_1 \\ ... \\ c_{n-1}
		\end{bmatrix} = c = M_{DFT, n} \cdot v = \frac{1}{n}
		\begin{bmatrix}
		1 & 1 & 1 & ... & 1 \\
		1 & \compl{\omega}^1 & \compl{\omega}^2 & ... & \compl{\omega}^{(n-1)}\\
		... & ... & ... & ... & ...\\
		1 & \compl{\omega}^{n-1} & \compl{\omega}^{2 \cdot (n-1)} & ... & \compl{\omega}^{(n-1)(n-1)}
		\end{bmatrix} \cdot 
		\begin{bmatrix}
		v_0 \\ v_1 \\ ... \\ v_{n-1}
		\end{bmatrix}$$ \n 
		Das Inverse: 
		$$\begin{bmatrix}
		v_0 \\ v_1 \\ ... \\ v_{n-1}
		\end{bmatrix} = c = M_{IDFT, n} \cdot c = 
		\begin{bmatrix}
		1 & 1 & 1 & ... & 1 \\
		1 & \omega^1 & \omega^2 & ... & \omega^{(n-1)}\\
		... & ... & ... & ... & ...\\
		1 & \omega^{n-1} & \omega^{2 \cdot (n-1)} & ... & \omega^{(n-1)(n-1)}
		\end{bmatrix} \cdot 
		\begin{bmatrix}
		c_0 \\ c_1 \\ ... \\ c_{n-1}
		\end{bmatrix}$$
	}
	\btask{
		
		\begin{enumerate}
			\item $DFT((0, 1)^T)$?
			\item $DFT((1, 1, 1, 0)^T)$?
			\item $DFT((1, 0, 3, 0)^T)$?
			\item $IDFT((0, 1, 0, 0)^T)$?
		\end{enumerate}
		
	}
	\bsol{
		
		\begin{enumerate}
			\item $DFT((0, 1)^T) =$ $$\frac{1}{2} \cdot (1, -1)^T$$
			\item $DFT((1, 1, 1, 0)^T) =$ $$\frac{1}{4} \cdot (3, -i, 1, i)^T$$
			\item $DFT((1, 0, 3, 0)^T) = $ $$\frac{1}{4} \cdot (4, -2, 4, -2)^T$$
			\item $IDFT((0, 1, 0, 0)^T) = $ $$(1, i, -1, -i)^T$$
		\end{enumerate}
	}

%----------------------------------------------------------------------------------
\bsubsec{Schnelle inverse Fourier Transformation}{
		
		Viel zuuu nervig, hier gibt es noch nichts zu sehen!
	}
	%\btask{
	%	
	%	TBD
	%	\begin{enumerate}
	%		\item TBD
	%	\end{enumerate}
	%	
	%}
	%\bsol{
	%	
	%	\begin{enumerate}
	%		\item TBD
	%	\end{enumerate}
	%}
%----------------------------------------------------------------------------------
\bquestion{
	}
	\btask{
		
		Beantworte folgende (Trick-)Fragen möglichst genau:
		\begin{enumerate}
			\item Was bedeutet es mathematisch, dass $DFT$ und $IDFT$ Umkehroperationen sind?
			\item Auf was für einem Prinzip beruht die Fourier-Transformation?
		\end{enumerate}
	}
	\bsol{
		
		\begin{enumerate}
			\item Es gilt: $IDFT(DFT(v)) = v$
			\item Interpolation mit Basisfunktionen
		\end{enumerate}
	}
%----------------------------------------------------------------------------------